\subsection{Beyond one tower: vision--language models}
\label{sec:exp-vlm}

Nothing in \cref{sec:method} is specific to a decoder-only stack. The risk is a KL divergence between
a frozen model and a masked copy of it; the estimator is a Gram matrix of single-unit ablation
features; the selector is a cost-normalised greedy under one budget. Point those three objects at a
vision--language model and the only thing that changes is that the unit inventory grows a second
tower. Everything else is held fixed: the same $M$ single-unit ablations, the same
$\Hmat_{uv}=\bar{D}_u^{\top}\bar{D}_v/P$, the same $\lambda$ protocol of \cref{sec:method-lambda}, the
same guard, no gradients, no fine-tuning and no adapter. One piece of bookkeeping is new --- the layer
index is composite, so the per-layer cap treats a vision layer and a language layer as distinct
layers rather than pooling them --- and it introduces no hyperparameter.

\textbf{Protocol.} The table below reports \vlmNcolsW\ vision--language models spanning $2$B to
$8$B and three model families --- Qwen2-VL \citep{wang2024qwen2vl}, Qwen2.5-VL
\citep{bai2025qwen25vl}, Qwen3-VL \citep{qwen3vl2025}, InternVL3 \citep{zhu2025internvl3} and
LLaVA-1.5 \citep{liu2024llava15}; the full grid of \vlmGridModelsW\ models at three ratios, adding
LLaVA-NeXT \citep{liu2024llavanext}, is in \cref{app:vlm}, on the
seven benchmarks this literature reports as a matter of course --- MMBench-EN
\citep{liu2024mmbench}, SEED-Bench \citep{li2024seedbench}, ScienceQA-IMG \citep{lu2022scienceqa},
AI2D \citep{kembhavi2016ai2d}, MMStar \citep{chen2024mmstar}, POPE \citep{li2023pope} and MME
\citep{fu2023mme} --- at $1000$ items each, with the dense model measured
under our own harness rather than quoted. Calibration is a fixed draw of Flickr30k image--caption pairs
\citep{young2014flickr30k}, whose images are disjoint from the COCO-derived suites by construction;
\cref{app:vlm} gives the draw for each model. The baselines are the
training-free structured selectors this field returns to: Magnitude \citep{han2015learning},
Wanda-sp \citep{sun2024wanda}, FLAP \citep{an2024flap}, LLM-Pruner's gradient Taylor score
\citep{ma2023llmpruner}, a ShortGPT-style layer drop \citep{men2024shortgpt}, and Random. Every arm receives the identical
budget, guard, calibration draw and scoring code; the only thing that differs between rows is which
units get removed. \Cref{app:vlm} gives the per-benchmark breakdown and every configuration measured.

\begin{table}[t]
\centering\small
\caption{Vision--language models, mean accuracy over seven benchmarks at $1000$ items each. Every
selector is training-free and runs under one shared budget and one shared guard. Retention is
floor-corrected against the dense model measured here.}
\label{tab:vlm}
% Five model groups plus the widest row label ("diagonal only") run 44pt past the text block at the
% default column separation. Tightening the gutter fits it without shrinking the type.
\setlength{\tabcolsep}{2pt}
\begin{tabular}{@{}lccccc@{}}
\toprule
 & \multicolumn{1}{c}{Qwen2-VL-2B} & \multicolumn{1}{c}{Qwen2.5-VL-7B} & \multicolumn{1}{c}{Qwen3-VL-8B} & \multicolumn{1}{c}{InternVL3-8B} & \multicolumn{1}{c}{LLaVA-1.5-7B} \\
\cmidrule(lr){2-2}\cmidrule(lr){3-3}\cmidrule(lr){4-4}\cmidrule(lr){5-5}\cmidrule(lr){6-6}
Selector & $20\%$ & $20\%$ & $20\%$ & $20\%$ & $20\%$ \\
\midrule
Dense (no pruning) & 73.0 & 81.6 & 83.4 & 83.6 & 64.1 \\
\midrule
\textbf{\method} & \textbf{59.3} & \textbf{74.4} & \textbf{71.1} & \textbf{75.9} & \textbf{57.1} \\
\quad diagonal only ($\lambda{=}0$) & 53.7 & 73.5 & 66.5 & 73.9 & 50.0 \\
\midrule
LLM-Pruner & 57.0 & 72.2 & 65.8 & 71.8 & 53.7 \\
Wanda-sp & 51.5 & 62.9 & 40.5 & 49.8 & 36.0 \\
FLAP & 36.6 & 73.3 & 68.4 & 74.8 & 41.1 \\
ShortGPT (layer drop) & 35.5 & 35.3 & 35.2 & 43.4 & 34.6 \\
Magnitude & 40.3 & 72.9 & 57.0 & 70.7 & 48.1 \\
Random & 48.0 & 64.4 & 53.8 & 67.9 & 52.8 \\
\midrule
\emph{\method retention of dense} & 66\% & 86\% & 76\% & 85\% & 78\% \\
\bottomrule
\end{tabular}

\end{table}

\textbf{Result.} \method leads every baseline in all \vlmNcolsW\ columns, by \vlmMarginLo\ to
\vlmMarginHi\ points of seven-benchmark mean, and is first outright in \vlmNfirst\ of the
\vlmNcells\ per-benchmark cells against a standard error of $\vlmSE$ points. No single baseline is the
one being beaten: the runner-up is \vlmRunnerUps. Because the budget is shared and the objective is
cost-normalised, the split between towers is an output rather than a hyperparameter: on the five
models whose vision tower holds under a tenth of the removable parameters it spends nothing there at
all, and on Qwen2-VL-2B, where the vision tower holds a third of them, it takes a fifth of the budget
from the vision side. Nothing in the objective was told to do either.

\textbf{Along the ratio.} \Cref{tab:vlm-sweep} carries Qwen2.5-VL-7B and Qwen3-VL-8B to $\rho{=}30\%$. The margin grows with the budget rather than shrinking --- $1.1$ to $4.9$ points on
Qwen2.5-VL-7B and $2.6$ to $10.8$ on Qwen3-VL-8B --- which is what a second-order objective should do
if the interactions it models are real: the more units leave together, the more the coupling between
them matters. \Cref{app:vlm} carries the same axis out to $\rho{=}40\%$ on all
\vlmModelsAtHighW\ models of the grid.

\textbf{Across the grid.} \Cref{app:vlm} tables all \vlmAllCells\ model--ratio cells --- every
model, every arm, every benchmark. At $\rho{=}20\%$ the result is not confined to the columns above:
on all \vlmGridModelsW\ models no baseline is ahead of \method by as much as $1$ point, and \method
is ahead by more than $2$ standard errors on \vlmNcolsW\ of them. At $\rho{=}40\%$ the axis has
closed --- no arm on any model retains a third of its dense score above chance --- and a
seven-benchmark mean no longer orders selectors that all sit within a few points of one another. What
separates methods is the budget in between, which is where \cref{tab:vlm-sweep} sits.

\begin{table}[t]
\centering\small
\caption{The ratio axis on Qwen2.5-VL-7B and Qwen3-VL-8B, $1000$ items per benchmark.
Same seven benchmarks, same arms, same guard as \cref{tab:vlm}.}
\label{tab:vlm-sweep}
\input{tables/tab_vlm_sweep}
\end{table}

\textbf{What the off-diagonal is worth.} Setting $\lambda{=}0$ keeps the budget, the guard and the
saliency term and deletes only the off-diagonal. \Cref{tab:vlm-ablation} runs that on every
vision--language model we built a curvature matrix for, at one ratio and one item count: it costs
between $\vlmAblLo$ and $\vlmAblHi$ points of seven-benchmark mean, $\vlmAblSigLo$ to
$\vlmAblSigHi$ times the standard error of the measurement, and it costs something on all
\vlmAblNW\ models. This is the same effect the text-only half of the paper reports, in an architecture
the estimator was never specialised for --- and because both arms are ours, run under one budget, one
guard and one calibration draw, it isolates the off-diagonal without routing the conclusion through
anyone else's implementation.

\begin{table}[t]
\centering\small
\caption{What the off-diagonal buys, per model, at $\rho{=}30\%$ and $1000$ items per benchmark.
$\lambda{=}0$ keeps the budget, the guard and the diagonal saliency term and removes only the
coupling; $\Delta$ is what that costs.}
\label{tab:vlm-ablation}
\begin{tabular}{@{}lrrrr@{}}
\toprule
Model & Dense & \method & \quad$\lambda{=}0$ & $\Delta$ \\
\midrule
Qwen3-VL-8B & 83.4 & \textbf{57.7} & 50.1 & $+7.6$ \\
LLaVA-1.5-7B & 64.1 & \textbf{43.6} & 36.4 & $+7.2$ \\
Qwen2.5-VL-3B & 77.4 & \textbf{50.0} & 42.9 & $+7.1$ \\
InternVL3-8B & 83.6 & \textbf{62.4} & 55.5 & $+6.9$ \\
Qwen2.5-VL-7B & 81.6 & \textbf{63.6} & 61.1 & $+2.5$ \\
LLaVA-NeXT-7B & 70.0 & \textbf{53.0} & 50.7 & $+2.3$ \\
Qwen2-VL-2B & 73.0 & \textbf{39.7} & 37.8 & $+1.9$ \\
\bottomrule
\end{tabular}

\end{table}
